\chapter{Literature Survey}
\label{chap:lit-survey}
\noindent \rule{6.6in}{0.01in}
\newpage

This chapter surveys prior work along three lines that together shape the design choices made in this thesis: machine-learning approaches to predicting the severity of a traffic incident, deep-learning approaches to forecasting traffic congestion more generally, and the ensemble/stacking techniques used to combine the four base classifiers in Chapter~\ref{chap:architecture}.

\section{Machine Learning for Incident and Crash Severity Prediction}

Predicting how severe a road-accident record is has attracted a fair amount of machine-learning attention, generally using post-hoc accident databases rather than real-time event reports. Dong et al.\ compared several boosting-based ensembles -- including LightGBM and other Gradient Boosting variants -- for road-traffic injury severity on highway accident data, and paired the best model with SHapley Additive exPlanations (SHAP) to identify which risk factors, such as month of year, accident cause, and driver age, drove its predictions most strongly, finding LightGBM to be the strongest of the boosting variants tried~\cite{dong2022predicting}. That finding lines up with what is observed independently in this thesis: among the four base classifiers evaluated in Chapter~\ref{chap:results}, LightGBM turns out to be the strongest single model at recovering the rarer Medium and Critical severity classes.

Iranitalab and Khattak took a broader comparative view, benchmarking logistic regression, nearest-neighbour classification, support vector machines, and neural networks against each other for crash-severity classification, and concluded that no single traditional model wins across every severity level -- an observation that motivates hybrid or ensemble strategies rather than a single best-in-class model~\cite{iranitalab2017comparison}. Zhu, Wang, and Li pushed in a different direction, proposing an ordinal-classification formulation that exploits the natural ordering of severity categories together with oversampling to counter class imbalance, and reported precision gains over standard nominal classifiers (neural networks, XGBoost, SVMs) once the imbalance was addressed explicitly~\cite{zhu2021crash}.

Feature selection and interpretability have also received attention in this space. Zhang et al.\ combined a Random Forest classifier with the Boruta wrapper algorithm to first narrow down the most informative accident attributes before training a final model, improving both interpretability and predictive performance on single- and multi-vehicle accident data~\cite{zhang2022hybrid}. Separately, Chen, Shao, and Ji paired CatBoost with SHAP explanations to study how driver experience and age jointly influence accident severity, which is one more instance of a broader pattern in this literature: gradient-boosted tree models paired with a post-hoc explainability tool such as SHAP~\cite{chen2021insights}.

\section{Deep Learning for Traffic Congestion Prediction}

A separate strand of work addresses the more classical problem of forecasting recurring traffic flow and congestion with deep learning rather than predicting the severity of a discrete incident. Yin et al.\ survey this space comprehensively, organising deep-learning traffic-prediction methods by the kind of spatio-temporal dependency they model -- grid-based, graph-based, or multivariate time-series~\cite{yin2022deep}. Kumar and Raubal narrow the focus specifically to congestion -- its detection, prediction, and alleviation -- and explicitly separate \emph{recurring} congestion, driven by regular commuting patterns, from \emph{non-recurring} congestion caused by discrete incidents, noting that the latter is comparatively under-served by the existing deep-learning congestion literature~\cite{kumar2021applications}. That observation is essentially the gap this thesis tries to sit in: the system built here is deliberately scoped to event-driven, non-recurring congestion rather than the recurring-flow problem that most of the deep-learning literature addresses. Graph- and recurrence-based architectures continue to be applied to the recurring side of the problem -- for instance, Hussain et al.\ use gated recurrent units over Internet-of-Vehicles traffic-flow data to model how flow evolves across junctions~\cite{hussain2023urban} -- but such approaches generally assume a continuous, real-time sensor stream that is not available in the event-log setting used in this thesis.

\section{Ensemble Learning and Stacked Generalization}

The core modelling idea used in Chapter~\ref{chap:architecture} -- training several classifiers independently and then learning, with a second model, how to combine their outputs -- follows the stacked-generalization framework introduced by Wolpert, in which the predictions of a set of base learners become the input features of a higher-level model~\cite{wolpert1992stacked}. Two of the four base learners used in this thesis are themselves well-established gradient-boosting frameworks: XGBoost, introduced by Chen and Guestrin as a scalable, regularised implementation of gradient boosting~\cite{chen2016xgboost}, and LightGBM, proposed by Ke et al.\ as a faster alternative that uses histogram-based, leaf-wise tree growth to scale more efficiently to large tabular datasets~\cite{ke2017lightgbm}. The fourth base learner, TabNet, was introduced by Arik and Pfister as a deep architecture built specifically for tabular data, using sequential attention to select a sparse subset of features at each decision step and thereby offering some built-in interpretability alongside its predictions~\cite{arik2021tabnet}. The multi-layer perceptron used as the third base learner is trained with the Adam optimizer, the adaptive gradient-based method introduced by Kingma and Ba that remains a default choice for training neural networks efficiently~\cite{kingma2015adam}.

\section{Handling Class Imbalance}

Severity datasets of this kind are essentially always imbalanced, since severe events are fortunately rarer than minor ones -- a pattern confirmed in this thesis's own dataset in Section~\ref{sec:target-construction}. The Synthetic Minority Over-sampling Technique (SMOTE), introduced by Chawla et al., addresses this by generating synthetic minority-class examples through interpolation between existing minority samples and their nearest neighbours, and remains one of the most widely used oversampling strategies for imbalanced classification more than two decades after its introduction~\cite{chawla2002smote}. Several of the crash-severity studies discussed above pair some form of imbalance handling -- oversampling, or merging adjacent severity categories -- with their classifiers specifically to improve minority-class detection~\cite{zhu2021crash}. This thesis takes a more modest first step in the same direction, using class-balanced sample weighting during training (Section~\ref{sec:base-classifiers}) and reporting per-class recall and macro-F1 alongside accuracy throughout Chapter~\ref{chap:results}, while leaving more elaborate imbalance-handling techniques such as SMOTE as a direction for future work (Section~\ref{sec:future-scope}).

\section{Summary and Positioning of this Thesis}

Taken together, the literature reviewed above supports three observations. First, gradient-boosted tree ensembles consistently perform well on structured, tabular incident-severity data, which is consistent with LightGBM and XGBoost being the two most reliable base models observed in this thesis. Second, deep-learning research on traffic congestion has concentrated mostly on the recurring, sensor-driven side of the problem, leaving event-driven, non-recurring congestion comparatively under-explored. Third, while stacked ensembles and imbalance-aware evaluation are each individually well established, they are not often applied together inside a single pipeline that also converts its prediction into an operational recommendation. This thesis sits at the intersection of these three threads: it applies a stacked ensemble of tree-based and neural classifiers to event-driven severity prediction, evaluates the result with explicit, front-loaded attention to class imbalance rather than accuracy alone, and -- unlike most of the severity-prediction work surveyed here -- couples the resulting prediction directly to a deterministic, rule-based resource-recommendation engine (Section~\ref{sec:recommendation-engine}) that turns a severity label into an actionable manpower and diversion plan.
